\usection{Tipo de Investigación}

Según su alcance, el presente trabajo de grado se enmarca dentro de la modalidad de investigación proyectiva. \citeA{hurtado2010} describe este tipo de investigación como una que “tiene como objetivo diseñar o crear propuestas dirigidas a resolver determinadas situaciones (...) Este tipo de investigación potencia el desarrollo tecnológico” (p.~133). Esta descripción encaja con el trabajo realizado dado que su objetivo fundamental es diseñar y desarrollar una propuesta de solución a una necesidad detectada por la carencia de herramientas de programación visual basadas en flujos de datos para el público infantil.

El proceso metodológico abarca desde el diagnóstico de los requerimientos educativos y técnicos, pasando por el diseño de la arquitectura y sintaxis del lenguaje, hasta la materialización de la propuesta en un entorno de software funcional validado por expertos en el área de educación.